\documentclass[11pt,letterpaper]{article}

\usepackage[margin=0.7in]{geometry}
\usepackage{amsmath, amssymb}
\usepackage{bm}
\usepackage{xcolor}
\usepackage{enumitem}
\setlength{\parindent}{0pt}
\setlength{\parskip}{0.4em}

\begin{document}

\begin{center}
  {\Large\bfseries Tumor--Immune Active Matter --- Key Equations}\\[0.2em]
  {\small I-Shan Tsai (\texttt{i3tsai@ucsd.edu}) \quad
          Chih-Yen Liu (\texttt{chl250@ucsd.edu})\quad UCSD, 2026-05-23}
\end{center}

\smallskip

%--------------------------------------------------------------------------
\textbf{1. T-cell chemotaxis} (the immune $\leftrightarrow$ tumor coupling --- defines the phase diagram)
\begin{equation*}
\frac{d\mathbf{R}_j}{dt}
  = v_I\,\mathbf{m}_j
  + \underbrace{\chi_a\,\nabla c_a(\mathbf{R}_j)}_{\text{attraction to tumor}}
  - \underbrace{\alpha\,\nabla\!\big[\,c_s + c_{IL10}\,\big](\mathbf{R}_j)}_{\text{repulsion by suppressant}}
  + \mathbf{F}^{\text{rep}}_j + \sqrt{2 D_I}\,\bm{\zeta}_j(t).
\end{equation*}
$\alpha \equiv \chi_s$ is the immunosuppression strength scanned across the phase diagram. Anti-PD-1 / anti-PD-L1 in clinic $\equiv$ lowering $\alpha$.

%--------------------------------------------------------------------------
\textbf{2. Two-scale reaction--diffusion fields} (the asymmetry that creates the phases)
\begin{equation*}
\frac{\partial c_a}{\partial t} = D_a\,\nabla^2 c_a + s_a\,\rho_T - \lambda_a\,c_a
\quad (\ell_a = \sqrt{D_a/\lambda_a} \approx 7\sigma:\ \text{long range})
\end{equation*}
\begin{equation*}
\frac{\partial c_s}{\partial t} = D_s\,\nabla^2 c_s + s_s\,\rho_T - \lambda_s\,c_s
\quad (\ell_s = \sqrt{D_s/\lambda_s} \approx 2\sigma:\ \text{short range})
\end{equation*}
$D_a \gg D_s$ is the defining asymmetry: T cells \emph{sense} the tumor from across the box but are \emph{excluded} only close in. The competition $\chi_a$ vs.\ $\alpha$ at fixed asymmetry produces clearance / control / escape.

%--------------------------------------------------------------------------
\textbf{3. Macrophage M1/M2 polarization} (the second treatment knob)
\begin{equation*}
\frac{d p_k}{dt} = \frac{p_{\text{eq}}(\mathbf{R}_k^{(M)}) - p_k}{\tau_p} + \sqrt{2 D_p}\,\zeta_k,
\qquad
p_{\text{eq}}(\mathbf{x}) = \tanh\!\big(\,b_{M_1}\,-\,\kappa_s\,c_s\,-\,\kappa_{IL}\,c_{IL10}\,\big)
\end{equation*}
\begin{equation*}
\frac{\partial c_{IL10}}{\partial t} = D_{IL10}\,\nabla^2 c_{IL10}
   + s_{IL10}\sum_k \max(0, -p_k)\,\delta(\mathbf{x} - \mathbf{R}_k^{(M)}) - \lambda_{IL10}\,c_{IL10}.
\end{equation*}
$b_{M_1}$ is the external M1-driving bias (CSF1R inhibitor / TLR agonist mimic). The $-\kappa_{IL}c_{IL10}$ term is the self-reinforcing IL-10 feedback that locks TAMs in M2 once the field is established.

%--------------------------------------------------------------------------
\textbf{4. EMT cell state} (the Model-2 cell-plasticity layer)
\begin{equation*}
\frac{d s_{T,i}}{dt} = \frac{s_{\text{eq}}(\mathbf{r}_i) - s_{T,i}}{\tau_s} + \sqrt{2 D_{\text{EMT}}}\,\eta_i,
\quad
s_{\text{eq}} = \mathrm{clip}\!\big(\,k_{\text{hyp},s}\,H + k_{\text{ECM},s}\,\rho_E + k_{\text{supp},s}\,c_s - k_{\text{MET}}\,,\,0,1\big).
\end{equation*}
Hypoxia indicator $H(\mathbf{x}) = \mathbb{1}[\,c_{O_2} < O^\star\,]$. Mesenchymal tumor cells lose cadherin, gain speed, and secrete more MMP via linear interpolation in $s_T$.

%--------------------------------------------------------------------------
\textbf{5. Hypoxia couples to immunity and angiogenesis}
\begin{equation*}
p_{\text{kill}}^{\text{eff}}(\mathbf{r}_i) = p_{\text{kill}}\,
   \big(1 - k_{\text{kill,hyp}}\,H(\mathbf{r}_i)\big)
\qquad
\frac{\partial c_{O_2}}{\partial t}
  = D_O\,\nabla^2 c_{O_2} - q_T\,\rho_T - \lambda_O\,c_{O_2}
  + s_{O_2}^{\text{ves}}\!\!\sum_v \delta(\mathbf{x}-\mathbf{x}_v)
\end{equation*}
\begin{equation*}
\frac{d\mathbf{x}_v}{dt} = \chi_V\,\nabla c_{\text{VEGF}}(\mathbf{x}_v)
\qquad
\text{vessel sprouts when } c_{\text{VEGF}}(\mathbf{x}_v) > c_V^\star.
\end{equation*}
Hypoxic regions are killed less efficiently \emph{and} secrete more suppressant; they also drive VEGF, which sprouts new vessels that restore $O_2$ --- closing the angiogenic loop.

\smallskip
\hrule
\smallskip

%--------------------------------------------------------------------------
\textbf{Order parameter} (the experimental observable)
\begin{equation*}
\Phi(\rho_I, \alpha) \;=\; \mathrm{clip}\!\left(\frac{N_T(T_f)}{N_T(0)},\,10^{-2},\,10^{+2}\right),
\qquad
\bar\Phi \;=\; \exp\!\Big(\,S^{-1}\textstyle\sum_s \log \Phi^{(s)}\Big).
\end{equation*}
Geometric mean over $S$ seeds. Sharp bistable transition in $(\rho_I, \alpha) \to$ \emph{clearance / control (dormancy) / escape}.

\vspace{0.4em}
\hrule
\vspace{0.4em}

{\small\itshape Full model + extensions + parameters: \texttt{docs/model.tex} (this folder). Code: \texttt{src/sim.py}, \texttt{sim\_extended.py}, \texttt{sim\_combined.py}, \texttt{sim\_tme.py}, \texttt{sim\_biophysical.py}.}

\end{document}
